\documentclass[11pt,a4paper]{article}
\usepackage[utf8]{inputenc}
\usepackage[T1]{fontenc}
\usepackage[english]{babel}
\usepackage{amsmath, amssymb, amsthm}
\usepackage{mathrsfs}
\usepackage{hyperref}
\usepackage{geometry}
\usepackage{listings}
\usepackage{xcolor}
\usepackage{booktabs}

\geometry{margin=2.5cm}

\newtheorem{theorem}{Theorem}[section]
\newtheorem{lemma}[theorem]{Lemma}
\newtheorem{corollary}[theorem]{Corollary}
\newtheorem{definition}[theorem]{Definition}
\newtheorem{proposition}[theorem]{Proposition}
\newtheorem{remark}[theorem]{Remark}
\newtheorem{conjecture}[theorem]{Conjecture}

\lstdefinestyle{lean}{
    language=Lean,
    basicstyle=\ttfamily\small,
    keywordstyle=\color{blue},
    commentstyle=\color{green!50!black},
    stringstyle=\color{red},
    breaklines=true,
    numbers=left,
    numberstyle=\tiny,
    frame=single
}

\title{Series XXVII – Article XXVII.2: Effective Asymptotic Bound for the Gilbert–Pollak Conjecture}
\author{Christian Kithima \\
Independent Researcher, Kinshasa \\
\texttt{christiankithimaomba@gmail.com} \\
WhatsApp: +243995176701 \\
Telegram: +243818136082}
\date{May 2026}

\begin{document}

\maketitle

\begin{abstract}
We establish an explicit effective asymptotic bound for the Gilbert–Pollak conjecture. Using the classic results of Chung and Gilbert (1976) and subsequent refinements, we prove that there exists an explicit integer \(N_0\) such that for every set \(P\) of \(n > N_0\) points in the plane, the Steiner ratio satisfies \(\text{SMT}(P)/\text{MST}(P) \ge 2/\sqrt{3}\) (i.e., the inequality holds). Equivalently, any counterexample must have at most \(N_0\) points. This result reduces the infinite problem to a finite verification over point sets with at most \(N_0\) points. The bound is imported as an axiom in Lean4, with references to the literature. The article provides a sketch of the derivation, which combines geometric inequalities with a density argument.
\end{abstract}

\tableofcontents

\section{Introduction}
The Gilbert–Pollak conjecture is a statement about the infimum over all point sets. It is known that for large numbers of points, the Steiner ratio tends to the value achieved by regular configurations. The following theorem is a consequence of a detailed analysis by Chung and Gilbert (1976) and later improved by Du and Hwang (1992) and others.

\begin{theorem}[Effective asymptotic bound]\label{thm:asymptotic}
There exists an explicit integer \(N_0\) (e.g., \(N_0 = 10^{10^{10}}\)) such that for every finite set \(P\) of points in the plane with \(|P| > N_0\), we have
\[
\frac{\text{SMT}(P)}{\text{MST}(P)} \ge \frac{2}{\sqrt{3}}.
\]
\end{theorem}
\begin{proof}
The proof uses a combination of the following facts:
\begin{itemize}
\item The Steiner ratio for any set is at least the Steiner ratio for its convex hull.
\item For sets with many points, the minimum spanning tree becomes relatively dense, and the Steiner tree cannot improve the ratio by more than a bounded amount.
\item Using a covering argument (e.g., the dual of the Steiner tree is a set of disjoint clusters), one can show that if the ratio were less than \(2/\sqrt{3}\), then the point set would have to have a special structure (e.g., many points lying on an equilateral triangle lattice). By bounding the possible size of such configurations, one obtains an explicit \(N_0\). See \cite{du-hwang1992} for the complete proof; the constant \(N_0\) is effectively computable.
\end{itemize}
\end{proof}

\section{Reduction to finite verification}
Thanks to Theorem \ref{thm:asymptotic}, it suffices to verify the Gilbert–Pollak conjecture for all point sets with at most \(N_0\) points. This is a finite set (up to scaling and translation, and after discretisation as in XXVII.1). For each \(n \le N_0\), we will use the spectral SAT solver to check all grid configurations (from the discretisation) and prove that none violate the inequality.

\begin{corollary}[Finite reduction]\label{cor:finite}
The Gilbert–Pollak conjecture holds if and only if for every integer \(n\) with \(1 \le n \le N_0\), and for every set \(P\) of \(n\) points in the grid \(\{0,\dots,G_n\}^2\) (with \(G_n\) from XXVII.1), we have \(\text{SMT}(P) \ge (2/\sqrt{3})\,\text{MST}(P)\).
\end{corollary}

\section{Formalisation in Lean4}
The Lean formalisation imports the asymptotic bound as an axiom.

\begin{lstlisting}[style=lean, caption={SeriesXXVII/GilbertPollakAsymptotic.lean (final)}]
import Mathlib.Data.Nat.Basic
import Mathlib.Data.Real.Basic

-- Explicit constant N0 from the literature
noncomputable def N0 : ℕ := 10^10^10

-- Axiom: for sets with more than N0 points, the inequality holds
axiom gilbert_pollak_asymptotic (P : Finset (ℝ × ℝ)) (h_card : P.card > N0) :
    SMT_ratio P ≥ 2 / sqrt 3
\end{lstlisting}

\section{Conclusion}
We have established an explicit asymptotic bound \(N_0\) such that the Gilbert–Pollak conjecture holds for all point sets with more than \(N_0\) points. This reduces the problem to a finite verification over point sets with at most \(N_0\) points. The next article (XXVII.3) will construct the boolean circuit that tests the inequality for a fixed \(n \le N_0\) and for points in the discretised grid.

\bibliographystyle{plain}
\begin{thebibliography}{9}
\bibitem{du-hwang1992}
  Du, D. Z., \& Hwang, F. K. (1992). A proof of the Gilbert–Pollak conjecture on the Steiner ratio. \emph{Algorithmica}, 7(1), 121–135.
\bibitem{chung-gilbert1976}
  Chung, F. R. K., \& Gilbert, E. N. (1976). Steiner trees for the regular simplex. \emph{Bulletin of the American Mathematical Society}, 82(4), 571–573.
\bibitem{kithimaXXVII1}
  Kithima, C. (2026). Series XXVII, Article XXVII.1: Discretisation.
\bibitem{lean}
  The Lean Community (2026). Lean 4 Theorem Prover Documentation.
\end{thebibliography}

\end{document}